% Print and mobile editions use the same source.
\ifdefined\mobile
\documentclass[aps,prd,reprint,onecolumn,nofootinbib,amsmath,amssymb]{revtex4-2}
\usepackage[paperwidth=270pt,paperheight=585pt,textwidth=246pt,textheight=545pt,top=18pt,left=12pt,headsep=7pt]{geometry}
\else
\documentclass[aps,prd,reprint,nofootinbib,amsmath,amssymb]{revtex4-2}
\fi
\usepackage[T1]{fontenc}
\usepackage{lmodern,amsthm}
\usepackage[hidelinks]{hyperref}
\hypersetup{pdfauthor={OpenAI Codex},pdftitle={Ten and twenty-six from volume-normalized curvature residues}}
\ifdefined\mobile\makeatletter\def\frontmatter@abstractwidth{246\p@}\makeatother\fi
\emergencystretch=1.5em
\newtheorem{theorem}{Theorem}
\newtheorem{proposition}{Proposition}
\newcommand{\Tr}{\operatorname{Tr}}
\newcommand{\Res}{\operatorname*{Res}}
\newcommand{\dv}{\,d\mathrm{vol}_g}
\begin{document}
\title{Ten and twenty-six from volume-normalized curvature residues}
\author{OpenAI Codex}
\date{27 November 2026}
\begin{abstract}
We give a geometric characterization of the dimensions in a unit-curvature
criterion proposed by Rivero and in its supersymmetric extension. On every
closed Riemannian manifold of dimension \(d>2\), the volume-normalized residue
of the Hodge Laplacian on all forms has curvature coefficient \(-(d-2)/24\).
Subtracting the scalar-sector residue gives \(-(d-2)/8\). Equality in
magnitude with the scalar curvature measured by small geodesic balls holds
at \(d=26\) and \(d=10\), respectively, whenever curvature is detectable.
On a spin manifold the same pair comes from the Dirac supercharge and its
covariant kinetic energy. The ratio of two unit-sphere areas supplies
\((d-2)/2\). We classify the remaining factor by the trace of the curvature
endomorphism and give a conformal scalar operator for which neither
denominator is universal. With absolute boundary conditions the same
coefficients multiply the sum of bulk scalar curvature and twice the
boundary mean curvature. The geometric equalities require no worldsheet.
Using either equality to select a physical dimension requires a matching
postulate.
\end{abstract}
\maketitle

\section{The coefficient fixed in 2004}

How does spin change the propagation of a particle on a curved manifold?
This question gives a way to vary the observable in the \(d=26\) proposal
while keeping its unit of volume. The number to be matched is the local
coefficient multiplying the scalar curvature in a specified residue.
We vary the dimension \(d>2\), fix the metric normalization of the principal
symbol, and use the same prescription in each dimension.

Four dimensions give a familiar geometric example. For a massless scalar
with conformal coupling, the operator is the positive scalar Laplacian
plus \(\xi_d R\), where \(\xi_d=(d-2)/[4(d-1)]\). Its first local heat
coefficient after the volume term is \cite{Physres1ScalarHeat}
\begin{equation}
 \left(\tfrac16-\xi_d\right)R
      =\frac{4-d}{12(d-1)}\,R .
 \label{eq:four}
\end{equation}
At \(d=4\) the geometric \(+R/6\) from the heat kernel and the
\(-\xi_4 R=-R/6\) from the conformal coupling cancel at every point,
even on a round sphere. The volume term survives, and higher coefficients
can still detect curvature. Four is distinguished here by a zero. The
two dimensions studied below are distinguished by equality with the
scalar curvature measured by small geodesic balls.

Let \(M^d\) be a closed Riemannian spin manifold. Write \(R\) for its scalar
curvature, positive on a round sphere, and \(D\) for the self-adjoint Dirac
operator of its Levi-Civita spin connection. On the classical triple
\((C^\infty(M),L^2(M,S),D)\), the normalized residue of
\(\lvert D\rvert^{2-d}\) gives the Einstein--Hilbert integral. The 2004 note
proposes that its dimensionless coefficient, with the Euclidean minus sign
taken out, equal one. The note also leaves open what quantization would
enforce that condition \cite{Rivero2004Critical,KalauWalze1995}.

Here is the normalization explicitly. Let \(P=\nabla^*\nabla+V\) be a
nonnegative Laplace-type operator on a Hermitian bundle of complex rank
\(r\), let \(f\) be a smooth test function, and put \(p=d/2-1\). Define
\begin{equation}
 \mathcal E_d(P;f)=
 \frac{(4\pi)^{d/2}\Gamma(d/2)}{r}
 \Res_{z=p}\Tr(fP^{-z}).
 \label{eq:residue}
\end{equation}
Complex powers are zero on the kernel. The identical prefactor at the volume
pole \(z=d/2\) gives \(\int_M f\dv\), for every \(P\) above. Equivalently,
the residue in Eq.~\eqref{eq:residue} is one half the Wodzicki residue of
\(fP^{-p}\), since \(P\) has order two. Inserting \(f\) tests the density
even on metrics with zero total scalar curvature.

Write \(Z_P(t;f)=(4\pi t)^{d/2}\Tr(f e^{-tP})/r\). The local heat
expansion through its first curvature term gives \cite{McKeanSinger1967,KalauWalze1995}
\begin{equation}
 \begin{aligned}
 Z_P(t;f)&=\int_M f\,[1+t b_P+O(t^2)]\dv ,\\
 b_P&=R/6-\operatorname{tr}(V)/r .
 \end{aligned}
 \label{eq:heat}
\end{equation}
Taking the Mellin residue multiplies \(b_P\)
by \(\Gamma(d/2)/\Gamma(p)=p\). For \(D^2\) the Lichnerowicz identity gives
\(V=R/4\). The integrated result is
\begin{equation}
 \mathcal E_d(D^2;f)=-\frac{d-2}{24}\int_M fR\dv .
 \label{eq:original}
\end{equation}
Unit matching gives \(d=26\). It fixes a dimension only when the coefficient
is matched before introducing an adjustable gravitational coupling.

\section{The spinning particle}

Take the worldline supercharge \(Q=D\) and normalize its Hamiltonian by
\(H=Q^2\), so the principal symbol is \(\lvert\xi\rvert^2\).
In even dimension \(Q\) is odd under spinor chirality. For a graded
realization in odd dimension one may double the spinors and put \(D\) in the
two off-diagonal blocks. Division by the fibre rank gives the same local
calculation. Set \(K=\nabla^{S*}\nabla^S\), the positive kinetic operator
built from that same connection. The supersymmetry algebra and the
Lichnerowicz identity give
\begin{equation}
 H=K+\tfrac14 R\,\mathbf 1_S .
 \label{eq:square}
\end{equation}
We propose to measure the gravitational response of this spin coupling by
subtracting the kinetic residue. This uses one bundle, one volume trace
and one Laplace symbol. \(K\) is a reference Hamiltonian. A second physical
particle with Hamiltonian \(K\) is unnecessary.

\begin{theorem}\label{thm:relative}
For the classical spin triple with the above connection and normalization,
\begin{equation}
 \mathcal E_d(H;f)-\mathcal E_d(K;f)
 =-\frac{d-2}{8}\int_M fR\dv .
 \label{eq:relative}
\end{equation}
If this relative coefficient must equal \(-1\) as an identity of local
densities, its unique dimension is ten.
\end{theorem}
\begin{proof}
The two heat traces have the same volume and the same \(R/6\) term.
Equation~\eqref{eq:square} leaves \(-R/4\) in their difference.
Multiplication by \(p\) gives Eq.~\eqref{eq:relative}.
\end{proof}

There is also a direct symbol calculation. The resolvent identity puts
one insertion of \(R/4\) between the resolvents of \(K\) and \(H\).
In the complex-power formula the leading symbol of the difference is
therefore the derivative of \(x^{-p}\) at \(x=\lvert\xi\rvert^2\) times
that insertion. Symbol derivatives and further insertions have lower order.
Thus \(H^{-p}-K^{-p}\) has order at most \(-d\), with symbol of degree
\(-d\) equal to \(-p R\lvert\xi\rvert^{-d}\mathbf 1_S/4\).
Integration on the unit cosphere in the volume normalization gives the same
answer for nonconstant \(R\).

At \(d=10\) the spin and kinetic coefficients are \(-1/3\) and \(+2/3\).
Their difference is \(-1\). At \(d=26\) they are \(-1\) and \(+2\).
These values specify which observable distinguishes the two dimensions.
If the reference is allowed its own curvature coupling
\(K_\xi=K+\xi R\), its subtraction gives
\(-p(1/4-\xi)\int fR\dv\). This follows locally from Eq.~\eqref{eq:heat}
and also holds on a round sphere for the nonnegative examples
\(\xi=0,1/6\). The latter choice returns the \(d=26\) coefficient.
The proposed reference is fixed by using covariant kinetic energy alone.

\section{Two traces of the de Rham particle}

The supersymmetric quantum mechanics of differential forms gives another
way to make the subtraction with a single Hamiltonian
\cite{WittenMorse1982}. Its complex supercharges are the exterior
derivative \(\mathrm d\) and its adjoint \(\mathrm d^*\).
The real supercharge \(Q_{\mathrm H}=\mathrm d+\mathrm d^*\) squares to
the Hodge Laplacian \(\Delta\). Form degree \(k\) counts occupied fermion
modes. The degree-zero sector contains all scalar wavefunctions, including
their excited eigenstates.

\begin{proposition}\label{prop:forms}
The heat trace of \(\Delta\) over all forms, divided by the fibre rank
\(2^d\), minus the scalar heat trace has vanishing leading volume term
and local curvature coefficient \(-R/4\). Its volume-normalized Mellin
residue therefore gives Eq.~\eqref{eq:relative}.
\end{proposition}
\begin{proof}
On \(k\)-forms write \(\Delta_k=\nabla^*\nabla+V_k\).
The fibre trace of \(V_k\) is an orthogonal invariant linear in curvature,
hence a multiple of \(R\). On a metric of constant sectional curvature
\(\kappa\), the Weitzenböck endomorphism is \(k(d-k)\kappa\,\mathbf 1\).
Each of the \(k\) indices pairs with \(d-k\) complementary directions in
the curvature action. Since \(R=d(d-1)\kappa\), for every metric
\begin{equation}
 \frac{\operatorname{tr} V_k}{\binom{d}{k}}
 =\frac{k(d-k)}{d(d-1)}\,R .
 \label{eq:degree}
\end{equation}
Under the ordinary trace each of the \(2^d\) occupation states has equal
weight. An ordered pair of distinct modes is occupied and empty,
respectively, with probability \(1/4\). Hence the mean of \(k(d-k)\) is
\(d(d-1)/4\). The all-form trace has coefficient \(R/6-R/4=-R/12\),
whereas \(k=0\) has \(R/6\).
\end{proof}

The dependence on the choice of trace can be computed. Weight each state
by \(y^k\), \(y>0\), and divide by \((1+y)^d\) to preserve the volume.
Each mode is occupied with probability \(y/(1+y)\). After subtracting
the scalar sector the magnitude of the curvature coefficient is
\begin{equation}
 C_d(y)=\frac{d-2}{2}\frac{y}{(1+y)^2}
 \ \leq\ \frac{d-2}{8}.
 \label{eq:fugacity}
\end{equation}
Within these positive factorized weights, invariance of the state weights
under Hodge duality \(k\leftrightarrow d-k\) fixes \(y=1\). Ten is the
lowest dimension in this family that can reach unit magnitude, and only
the ordinary trace reaches it at ten.

The supersymmetric index measures a different quantity. In the alternating
trace \(\Tr((-1)^k e^{-t\Delta})\), \(Q_{\mathrm H}\) pairs the two
parities at every positive eigenvalue. The global trace equals the Euler
characteristic for all \(t\) \cite{WittenMorse1982}. It supplies no
integrated curvature residue at \(p>0\). The degree-zero projector used in
Proposition~\ref{prop:forms} commutes with \(\Delta\) and generally fails
to commute with \(Q_{\mathrm H}\). The postulate specifies which spectral
sectors the gravitational observable compares.

\section{A volume comparison in every dimension}\label{sec:volume}

Both factors in the calculation have a geometric meaning. Write
\(\omega_d=\mathrm{Vol}(B^d_1)\) and \(s_j=\mathrm{Vol}(S^j_1)\)
for the Euclidean unit ball and sphere volumes. The elementary recursion
\(\omega_d=(2\pi/d)\omega_{d-2}\), or the Gaussian integral for \(s_j\),
gives
\begin{equation}
 \pi\,\frac{s_{d-3}}{s_{d-1}}
 =\frac{\Gamma(d/2)}{\Gamma((d-2)/2)}
 =\frac{d-2}{2}.
 \label{eq:areas}
\end{equation}
It is the ratio that converts the first curvature heat coefficient to a
residue when the leading residue measures Riemannian volume. For the
curvature itself there is a measurement using only distances and volumes.
At any \(x\in M\) set
\begin{equation}
 R_{\rm ball}(x)=
 \lim_{b\downarrow0}\frac{6(d+2)}{b^2}
 \left(1-\frac{\mathrm{Vol}_g(B_g(x,b))}{\omega_d b^d}\right).
 \label{eq:geodesic}
\end{equation}
This limit equals \(R(x)\) \cite{GrayVanhecke1979}. Indeed, in normal
coordinates the volume density is
\(1-\mathrm{Ric}_{ij}(x)u^iu^j/6+O(|u|^3)\).
The average of \(u^iu^j\) on a Euclidean ball of radius \(b\) is
\(\delta_{ij}b^2/(d+2)\). These two volume formulas fix the comparison
without choosing a length scale.

\begin{theorem}\label{thm:geometric}
Let \(M^d\) be any smooth closed Riemannian manifold, \(d>2\).
Let \(\Delta_{\rm all}\) be the Hodge Laplacian on the full complex
exterior algebra and \(\Delta_0\) the scalar Laplacian. Use
Eq.~\eqref{eq:residue} on each bundle with its own rank. As identities
for every smooth \(f\),
\begin{equation}
 \begin{aligned}
 \mathcal E_d(\Delta_{\rm all};f)
    &=-\frac{d-2}{24}\int_M fR_{\rm ball}\dv ,\\
 \mathcal E_d(\Delta_{\rm all};f)-\mathcal E_d(\Delta_0;f)
    &=-\frac{d-2}{8}\int_M fR_{\rm ball}\dv .
 \end{aligned}
 \label{eq:ballidentities}
\end{equation}
If \(R\) is nonzero somewhere, equality of the negative of the first
functional with \(\int fR_{\rm ball}\dv\) for all \(f\) holds exactly at
\(d=26\). For the second functional it holds exactly at \(d=10\).
The statements apply to nonorientable manifolds with the Riemannian
volume density and the \(L^2\) adjoint of the exterior derivative.
\end{theorem}
\begin{proof}
The orthogonal trace calculation in Proposition~\ref{prop:forms} is local
and uses the exterior algebra of each tangent space. It gives \(-R/12\)
and \(-R/4\) for the two heat densities. Multiply by the area ratio in
Eq.~\eqref{eq:areas} and insert Eq.~\eqref{eq:geodesic}. A nonnegative
test function supported where \(R\) has a strict sign then proves both
necessity statements.
\end{proof}

This is one precise sense in which the dimensions are critical in pure
geometry. The positive ratio of the negative spectral density to the
ball curvature, wherever \(R\ne0\), crosses one with slope \(1/24\)
or \(1/8\) in the two dimension formulas. At the crossing two independently
defined curvature measurements agree pointwise. Scalar-flat metrics give
zero in all three measurements, and a single integral with \(f=1\) may
lose the signal through sign cancellation. The comparison concerns the
usual scalar curvature and the residue calibrated at the volume pole.
Allowing an arbitrary prefactor for either observable would move the
crossing. In \(d=2\) the power \(P^0\) is the identity up to a smoothing
projection and has zero Wodzicki residue. In \(d=1\) scalar curvature
vanishes identically.

\subsection{Which geometric operators have these denominators?}

There is a necessary and sufficient local test. For any Laplace-type
operator \(P=\nabla^*\nabla+V\) put \(v=\operatorname{tr}V/r\).
Its ordinary residue density and its density relative to
\(\nabla^*\nabla\), with respect to Riemannian volume, are
\begin{equation}
 e_{\rm ord}=p(R/6-v),\qquad e_{\rm rel}=-p\,v .
 \label{eq:class}
\end{equation}
Equation~\eqref{eq:heat} proves both identities. If \(P\) is only bounded
below, discarding its finitely many nonpositive eigenvalues leaves these
positive-pole residues unchanged. For a specified positive \(k\), the
ordinary density is \(-(d-2)R/k\) everywhere exactly when
\(v=(1/6+2/k)R\). The relative density takes that value exactly when
\(v=2R/k\). These tests include points where \(R=0\). If only magnitudes
are compared, one takes the absolute values in Eq.~\eqref{eq:class}.
In particular the signed pair \(k=24\) ordinary and \(k=8\) relative
is equivalent to the single condition \(v=R/4\).

Every full spin Dirac operator coupled to a unitary auxiliary connection
by Clifford multiplication passes this test. In its square the additional
endomorphism is a sum of \(c(e_i)c(e_j)\) times the auxiliary curvature
for \(i<j\). Each Clifford product has zero spin trace, so the rank mean
of the endomorphism remains \(R/4\). The same argument holds locally
for the full spin\(^{c}\) Dirac operator with its determinant-line
connection. It thus applies, for example, to the canonical spin\(^{c}\)
operator \(\sqrt{2}(\bar\partial+\bar\partial^*)\) on complex projective
space with its Kähler metric in every complex dimension \(m\), with
\(d=2m\). The real dimensions ten and twenty-six correspond to \(m=5\)
and \(m=13\). The all-form version of Theorem~\ref{thm:geometric}
requires neither a spin nor a spin\(^{c}\) structure.

Other canonical choices on the very same metric answer the family
question differently. A minimal scalar has \(v=0\), ordinary coefficient
\(+(d-2)/12\), and zero relative coefficient. The conformal scalar
Laplacian has \(v=(d-2)R/[4(d-1)]\). Substitution gives
\begin{equation}
 \begin{aligned}
 e_{\rm ord}^{\rm conf}
  &= -\frac{(d-2)(d-4)}{24(d-1)}R ,\\
 e_{\rm rel}^{\rm conf}
  &= -\frac{(d-2)^2}{8(d-1)}R .
 \end{aligned}
 \label{eq:conformal}
\end{equation}
Neither is a dimension-independent \(1/8\) or \(1/24\) times
\(-(d-2)R\). The ordinary expression carries the zero of Eq.~\eqref{eq:four}
through volume normalization.
Projection to a fixed form degree provides another geometric choice with
\(v\) given by Eq.~\eqref{eq:degree}. Thus membership in a family of
manifolds alone cannot decide which denominator a probe measures.

\subsection{Families available in every real dimension}

For round \(S^d_b\) and for a closed hyperbolic quotient of curvature
\(-1/b^2\), the ball curvatures are \(+d(d-1)/b^2\) and
\(-d(d-1)/b^2\), respectively. Closed hyperbolic examples exist in every
dimension under consideration \cite{Millson1976}. The two signs give
opposite measured densities and the same two ratios. Real projective
\(\mathbb{RP}^d_b\), with the antipodal quotient metric, has the same
local densities as the sphere and half its integrated residues. This
also tests the bundle hypothesis. With \(a\) the degree-one mod-two
generator, \(T\mathbb{RP}^d\oplus\mathbf 1=(d+1)\gamma\) for the
tautological line \(\gamma\), so \(w(T)=(1+a)^{d+1}\). For \(d>2\)
there are ordinary spinors only for \(d=3\bmod 4\). At both ten and
twenty-six the form residues still realize the corresponding equalities.

An algebraic group family holds curvature fixed as dimension varies.
Take \(G_d=SU(2)\times T^{d-3}\) with the round bi-invariant metric of
radius \(b\) on \(SU(2)\) and flat torus factors. Then
\(R=6/b^2\) for every \(d\ge3\). More generally, multiplying any
closed seed with detectable scalar curvature by flat circles leaves its
pointwise ball curvature fixed. For tests pulled back from the seed,
divide integrated quantities by the flat volume. Adding one circle
increases the two dimensionless coefficients in magnitude by exactly
\(1/24\) and \(1/8\). A family starting at seed dimension \(m\) realizes
each crossing for which the corresponding dimension is at least \(m\).

Flat tori give zero readings in all dimensions. One can keep the topology
and vary the metric. For \(g=e^{2u}g_0\) on \(T^d\) with \(g_0\) flat,
integration of the conformal scalar-curvature formula by parts gives
\begin{equation}
 \int_{T^d} R_g\dv
 =(d-1)(d-2)\int_{T^d}e^{(d-2)u}
           |\nabla_0 u|^2\,d\mathrm{vol}_0 .
 \label{eq:torus}
\end{equation}
Every smooth nonconstant \(u\) makes this quantity positive and exposes
both coefficients on the same torus. Nonflat scalar-flat examples also
occur. A round two-sphere times a closed hyperbolic surface of the
opposite sectional curvature, times \(T^{d-4}\), has \(R=0\) for all
\(d\ge4\) although its sectional curvatures are nonzero. These metrics
give a test of exactly what the first curvature residue detects.

\subsection{Compact domains and the volume of an inward collar}

A smooth boundary contributes at the curvature pole itself. Let \(h\)
be the sum of principal curvatures for the outward normal, positive on
a Euclidean sphere. Write \(A(s)\) for the area of the hypersurface a
distance \(s\) inward from the boundary. The normal Jacobian has derivative
\(-h\) at zero, so \(A'(0)=-\int_{\partial M}h\,dA\). The quantity
\begin{equation}
 \mathcal R(M)=\int_M R_{\rm ball}\dv-2A'(0)
 \label{eq:collar}
\end{equation}
is measurable by interior balls and collar volumes, since the area
\(A(s)\) is the derivative of the collar volume. Use Eq.~\eqref{eq:residue}
with \(f=1\) and the heat spectrum of the indicated boundary problem.

\begin{proposition}\label{prop:boundary}
On any compact smooth Riemannian \(d\)-manifold with smooth boundary,
\(d>2\), give differential forms absolute boundary conditions and the
degree-zero sector Neumann conditions. The ordinary all-form residue
and its difference from the scalar residue are, respectively,
\(-(d-2)\mathcal R(M)/24\) and \(-(d-2)\mathcal R(M)/8\).
If \(\mathcal R(M)\ne0\), the same two unit crossings follow.
\end{proposition}
\begin{proof}
For the full heat trace the normalized boundary part of the first
curvature coefficient is \(h/3+2\operatorname{tr}(S)/r\), where
\((\nabla_{\rm in}+S)\omega=0\) on the Robin block and \(S\) is
extended by zero on the Dirichlet block. This is the mixed heat
coefficient with constant smearing \cite{BransonMixed1999}.
The absolute conditions are \(\iota_\nu\omega=\iota_\nu\mathrm d\omega=0\).
Components containing the normal covector vanish at the boundary. On a
tangential form with occupied principal directions \(I\), the second
condition gives \(S=-\sum_{i\in I}\kappa_i\) with the inward derivative.
Indeed, differentiating the normal covector in a principal tangent
direction supplies that direction's principal curvature in the condition
with the outward derivative. Of all \(2^d\) form components, one quarter
have a given tangential mode occupied and the normal mode empty.
Thus \(\operatorname{tr}(S)/2^d=-h/4\), and the full boundary density
is \(-h/6\). The Neumann scalar has \(S=0\) and density \(h/3\).
Together with the interior densities, the full and relative heat
coefficients are \(-\mathcal R(M)/12\) and \(-\mathcal R(M)/4\), in the
common heat factor. Equation~\eqref{eq:areas} finishes the proof.
\end{proof}

The Euclidean ball of radius \(b\) gives a family in every dimension
with zero interior curvature and
\(\mathcal R(M)=2(d-1)\mathrm{Area}(\partial M)/b\). Its entire signal
comes from the collar. Other Robin data contribute through their trace.
The bulk-to-boundary ratio two also occurs for the Dirac boundary
problem of Chamseddine and Connes \cite{ChamseddineConnesBoundary2007}.

\section{Checks on curved products}

An exact heat trace checks the sign and the normalization without using a
local heat coefficient for forms. Take \(M=S^2_b\times T^{d-2}\) with a
round sphere of radius \(b\) and a flat torus. Write \(Z_0^S(t)\) and
\(Z_0^T(t)\) for their scalar heat traces including zero modes.
Every nonconstant scalar spherical harmonic \(u\) gives the one-forms
\(\mathrm du\) and \(*\mathrm du\) at the same eigenvalue. They span the
one-form spectrum by Hodge decomposition and \(H^1(S^2)=0\).
Hodge duality identifies the two-form and scalar spectra. The sum over
sphere form degrees is consequently \(4Z_0^S-2\). Each torus form
component has the scalar spectrum, so factorization gives, for every
\(t>0\),
\begin{equation}
 2^{-d}\Tr(e^{-t\Delta_M})-\Tr(e^{-t\Delta_{0,M}})
 =-\tfrac12 Z_0^T(t).
 \label{eq:sphere}
\end{equation}
The right side starts with \(-\mathrm{Vol}(T^{d-2})/
[2(4\pi t)^{(d-2)/2}]\).
Since \(\int_{S^2_b}R\dv=8\pi\), this is precisely the \(-R/4\) term
in Eq.~\eqref{eq:heat}, independently of \(b\).

For a ten-to-nine comparison take \(Y=T^4\times S^5_B\) and
\(M=Y\times S^1_\rho\), where \(B\) and \(\rho\) are radii. This is a
smooth circle version of the round-sphere and interval example in the
Kaluza--Klein construction \cite{TwoEndpoint2026}. Use the periodic spin
structure on the circle and a test function pulled back from \(Y\).
Both \(H\) and \(K\) split into the corresponding operator on \(Y\) plus
the circle Laplacian. The restriction of the ten-dimensional spin
bundle consists of two copies of the nine-dimensional spinors, so division
by the respective ranks gives the same circle factor for both traces,
\begin{equation}
 \Theta_\rho(t)=\sum_{m\in\mathbb Z}e^{-tm^2/\rho^2}.
\end{equation}
For fixed \(\rho>0\), Poisson summation gives
\(\Theta_\rho(t)\sim 2\pi\rho/\sqrt{4\pi t}\) as \(t\to0\).
The ten-dimensional relative residue has density \(-R\), with
\(R=20/B^2\). Taking the circle zero mode first gives
\(\Theta_\rho\to1\) as \(t/\rho^2\to\infty\). The resulting
nine-dimensional operator has its own residue at \(p=7/2\), with density
\(-7R/8\). Thus the unit condition applies to the resolved smooth parent
in this example. Selecting a circle radius or a compactification path
requires dynamics beyond the local coefficient.

\section{Relation to the superstring}

There is a published string calculation of the same local curvature
difference. In the spectral action of Chamseddine, Eqs.~(50)--(52)
give equal leading volume terms for the Neveu--Schwarz and Ramond zero-mode
Hamiltonians. With three-form torsion set to zero, their next densities
carry \(+R/6\) and \(-R/12\), in the common factor used there. Equation
(53) combines them in the NS-minus-R orientation and obtains \(+R/4\)
\cite{ChamseddineSuperstring1997}. That calculation has already set the
target dimension to ten before the heat expansion. Theorem~\ref{thm:relative}
starts with variable \(d\), uses the opposite orientation to retain the
Euclidean sign of Eq.~\eqref{eq:original}, and applies the unit criterion.
Agreement concerns the curvature term. The full string action also carries
oscillator, projection and ghost contributions.

For comparison the flat Ramond--Neveu--Schwarz string with no extra matter
has, per chiral sector, one real scalar and one real Majorana fermion per
target direction. Their central charges are \(1\) and \(1/2\). The
reparametrization ghosts contribute \(-26\) and the commuting
superconformal ghosts contribute \(+11\). Cancellation of the conformal
anomaly therefore gives \cite{FriedanMartinecShenker1986}
\begin{equation}
 c_{\mathrm{tot}}=\tfrac32 d-26+11=0,\qquad d=10 .
\end{equation}
In the effective-string calculation the factor \((d-2)/24\) also occurs in
the vacuum energy of transverse oscillators on a cylinder
\cite{QCDStringCasimir}. That observable is an energy shift. In the
critical string, anomaly cancellation supplies the selection rule. In the
commutative spectral triple, supersymmetry fixes \(R/4\), volume calibration
fixes the normalization of the residue, and unit matching for the
\emph{relative} action supplies the remaining premise. A derivation from a
quantized parent would have to fix the subtraction and its gravitational
normalization independently. An allowed finite \(\int R\dv\) counterterm
with an unfixed coefficient would leave that matching adjustable.

\bibliography{towerless_intervals}
\end{document}
